\chapter{The P\'olya density as a maximum order statistic and the Thorin measure of the clock}\label{chap:maxwald}\chaptermark{The P\'olya density as a maximum}

This chapter follows the first author's manuscript \cite{P26mw}, which starts from one observation: the law behind the P\'olya kernel of the Riemann Xi function is built from the maximum $M_4$ of independent exponential random variables of Section~\ref{sec:maxexp:products}, in the sense that $e^{2Y}$ is a sum of two independent copies of $M_4$ (Proposition~\ref{mw:prop:sum}). Everything else follows from asking which parts of that observation survive the passage from the P\'olya density to the Thorin measure of the clock, where the Riemann Hypothesis lives.

The Thorin programme of this book asks whether the clock, with Laplace transform $1/G(\lambda)=\xi(\frac12)/\xi(\frac12+\sqrt\lambda)$, is a generalised gamma convolution. If it is, the Thorin measure is $\sum_\gamma\delta_{\gamma^2}$ and RH holds, and the two statements are equivalent (Proposition~\ref{prop:rigid}). The unconditional statements about the clock one level below the GGC property, its existence, infinite divisibility and self decomposability and the ladder of signs of the derivatives of its heat trace, are those of Section~\ref{ch:heat}, and this chapter uses them as they stand. What it adds to the book is the following, each item with its label and, where the book has the neighbouring result, the label of that result.

\begin{enumerate}
\item[(a)] The P\'olya kernel through the density of $e^Y=\sqrt{M_4+M_4'}$, with $\frac\pi2S_2$ a sum of two independent copies of $M_4$ (Proposition~\ref{mw:prop:sum}). Biane, Pitman and Yor \cite{BPY01} have $e^{2Y}=\frac\pi2S_2$, and Section~\ref{ch:clock} has the tilt that turns $Y$ into the P\'olya variable; the reading of $S_2$ as a sum of two maxima of independent exponential variables is the addition.
\item[(b)] The Khinchin factorisation of the tilted maximum: for $\sigma>1$, $\widetilde M_4\,4^{-K}$ has the law of $\Gamma_{1+\sigma/2}/(\pi N_\sigma^2)$, where $N_\sigma$ is the zeta distributed integer and $4^{-K}$ is an independent factor that modifies only the prime $2$ (Theorem~\ref{mw:thm:khinchin}). The zeta distribution is that of Sections~\ref{rz:sec:strip} and \ref{rz:sec:levystrip}, and the tilt is the size biasing of Section~\ref{bp:sec:mellin}.
\item[(c)] The language of Wald couples, and the logistic law as a normal scale mixture with Kolmogorov mixing, read through the maximum (Theorem~\ref{mw:thm:wald}). The pair $(\sin\theta/\theta,\theta/\sinh\theta)$ is the classical van Dantzig pair of Section~\ref{ch:clock} \cite{Lukacs68}, the logistic statement is due to Andrews and Mallows \cite{AndrewsMallows74} and Stefanski \cite{Stefanski91}, and the language is that of \cite{P18}. The clock forms a Wald couple with the P\'olya variable by Theorem~\ref{thm:ch12:idsd} (Remark~\ref{mw:rem:waldT}).
\item[(d)] A criterion for the Thorin condition in the variable $-x-i\eta$: the clock is a GGC if and only if the smoothed Thorin density $r_\eta$ is nonnegative, and an off line zero $\frac12+\delta+i\gamma$ shows in $r_\eta$ as an atom of negative mass at the explicit point $(x,\eta)=(\gamma^2-\delta^2,2\delta\gamma)$ (Theorem~\ref{mw:thm:stieltjes}, Lemma~\ref{mw:lem:detect}). The equivalences of RH with the GGC property and with the complete monotonicity of the heat trace are in the book (Proposition~\ref{prop:rigid} and Section~\ref{ch:clock}), the smoothing is that of \eqref{eq:ch3:inversion}, and the positivity criterion and the off line signature in the variable $\theta=\sqrt z$ are Remark~\ref{rem:ch2:oneline} and Proposition~\ref{prop:signature}; the criterion in the variable $-x-i\eta$ and the closed form of the detection are the addition.
\item[(e)] An explicit GGC law within a stated Laplace distance of the clock, built from the verified zeros (Proposition~\ref{mw:prop:approx}); the book has the GGC approximants $1/G_N$ from the verified zeros (Section~\ref{sec:weak:limits}), and the explicit distance is the addition.
\item[(f)] The exact form, in the variable $\lambda$, of the term of one zero in the Thorin density of the basepoint family (Proposition~\ref{mw:prop:lambda}), next to the Poisson kernels in the variable $\theta=\sqrt z$ of Proposition~\ref{prop:struct}.
\item[(g)] The clock is not a maximum of independent exponential variables (Proposition~\ref{mw:prop:nomax}): the maximum structure belongs to the theta function and not to the zeros. Section~\ref{sec:maxexp:products} ends with the remark that the construction stops at $\zeta$, and the residues of the clock are in Section~\ref{ch:clocklaw}.
\item[(h)] The Gumbel norming constant of the clock from Gnedenko's extreme value theory (Proposition~\ref{mw:prop:gnedenko}); the hazard rate of Section~\ref{ch:clocklaw} tends to $\gamma_1^2$, the reciprocal of that constant, and the reading through extreme value theory is the addition.
\item[(i)] The finite stages $\xi_N$ of the maximum, their symmetrisations $\xi_N^{\mathrm{sym}}$ and the symmetric stages $F_N^{\mathrm{sym}}$, with the heights up to which their zeros lie on the line (Theorem~\ref{mw:thm:finite}, Corollary~\ref{mw:cor:hurwitz}, Subsection~\ref{mw:sec:finiteN}), next to the truncations of the Biane, Pitman and Yor lattice of Section~\ref{rz:sec:finite} (Theorem~\ref{rz:thm:hurwitzrh}, Proposition~\ref{rz:prop:symtrunc}, Theorem~\ref{rz:thm:symhurwitz}).
\item[(j)] Hyperbolic complete monotonicity of the clock to a finite order from the Hadamard product, and the equivalence of RH with membership of the clock in the Bondesson class $\mathrm{BO}$ (Theorem~\ref{mw:thm:bondesson}), next to the ladder of Section~\ref{ch:heat} (Theorems~\ref{thm:ch12:cm} and \ref{thm:ch12:window}) and the HCM form of the Thorin condition (Corollary~\ref{cor:hcm}).
\item[(k)] A finite form of the basepoint density below $\alpha=1$, with explicit pole terms and a Fej\'er taper at a cut off $X$, whose limit as $X\to\infty$ needs every zero within $\alpha-\frac12$ of the line (Proposition~\ref{mw:prop:AX}), next to Theorem~\ref{ri:thm:ggcriesz}.
\end{enumerate}

\paragraph{What is not claimed.} Nothing here proves RH. Every statement equivalent to RH is labelled as such, and every unconditional statement is proved by direct computation, from the explicit formula, from finite verification, or from the pairing structure of the zeros, in this chapter or in the sections it cites.

\paragraph{The dictionary.} Table~\ref{mw:tab:dict} sets the trigonometric side, where everything is classical, against the Riemann side; the entries are proved or cited in the subsections that follow.

\begin{table}[htbp]
\centering\small
\begin{tabular}{p{0.25\textwidth}p{0.30\textwidth}p{0.35\textwidth}}
\toprule
 & Trigonometric side & Riemann side\\
\midrule
entire function & $\sin(\sqrt\pi\theta)/(\sqrt\pi\theta)$ & $\Xi(\theta)/\Xi(0)$\\
real zeros & $\theta=\pm\sqrt\pi\,n$ & \raggedright $\theta=\pm\gamma_k$ (all real iff RH)\tabularnewline
space component & \raggedright $\sqrt{\pi/2}\,\upsilon$, with $\upsilon$ uniform on $(-1,1)$ & \raggedright $X/\sqrt2$, with $X$ the P\'olya variable\tabularnewline
time component & $M_4$ & $T$, the clock\\
Laplace transform & \raggedright $\sqrt{\pi\lambda}/\sinh\sqrt{\pi\lambda}$ & \raggedright $1/G(\lambda)=\xi(\frac12)/\xi(\frac12+\sqrt\lambda)$\tabularnewline
exists as a law & yes & yes (Theorem~\ref{thm:ch12:idsd})\\
ID and SD & yes & yes (Theorem~\ref{thm:ch12:idsd})\\
\raggedright $x$ times the L\'evy density & $\sum_{n\ge1}e^{-\pi n^2x}$ & $W(x)=\sum_ke^{-a_kx}$\\
GGC, Thorin measure & $\sum_{n\ge1}\delta_{\pi n^2}$ & \raggedright $\sum_k\delta_{\gamma_k^2}$ if and only if RH (Proposition~\ref{prop:rigid})\tabularnewline
\raggedright maximum of exponentials & \raggedright yes, with rates $\pi n$ and $\pi(2m-1)$ & no (Proposition~\ref{mw:prop:nomax})\\
Gnedenko norming & $a_n=1/\pi$ & $a_n=1/\gamma_1^2$\\
\bottomrule
\end{tabular}
\caption{The dictionary between the trigonometric side, where everything is classical, and the Riemann side developed in this chapter. The variable $\theta$ is that of Section~\ref{ch:clock}; ID and SD abbreviate infinitely divisible and self decomposable, and the Wald couple of the third and fourth rows is that of Theorem~\ref{mw:thm:wald} and Remark~\ref{mw:rem:waldT}.}
\label{mw:tab:dict}
\end{table}

A reader interested in the probabilistic picture can start with Section~\ref{mw:sec:max}, and a reader interested in the Thorin side with Section~\ref{mw:sec:stieltjes}.
